\section{Reasoning Rerankers and Reinforcement Learning for Ranking}
\label{sec:rerank-rl}

A second-stage \emph{reranker} re-scores a small candidate set returned by a
first-stage retriever. With large language models (LLMs) this re-scoring can be
framed as a generation problem, and a 2024--2025 wave of work has equipped these
rerankers with explicit \emph{reasoning} and trained them with reinforcement
learning (RL). This section surveys the prompting paradigms, the reasoning
rerankers, the RL recipes that make reasoning feasible at small scale, and the
distillation techniques that internalise reasoning into a sub-1B (or even
non-generative) scorer. Throughout, the target application is a
\emph{causal-precedence} scorer of $<\!1$B parameters that judges whether one
procedural step enables another, reranks candidate follow-ups, and feeds a
workflow orderer (Sec.~\ref{sec:ordering}).

\subsection{LLM reranking paradigms}
\label{sec:rerank-paradigms}
The standard taxonomy distinguishes four prompting paradigms by \emph{how many
candidates enter a single prompt} and what the model emits, trading cost against
effectiveness \citep{rl:zhou2025rerankingsurvey}.

\paragraph{Pointwise.} One query--document pair per call; the model emits an
independent relevance score (via a yes/no token probability or a graded label).
Cost is $O(N)$ for $N$ candidates and the scores are \emph{directly comparable
and reusable}, but they require calibration (Sec.~\ref{sec:calibration}) and the
per-item signal is the weakest of the four.

\paragraph{Pairwise.} Two documents per call (``which is more relevant?''),
aggregated into an order by sorting. This yields the strongest relative
judgments but costs up to $O(N^2)$ comparisons, reduced to $O(N\log N)$ with
sorting networks. Pairwise prompting maps naturally onto \emph{causal precedence}
(``does step $A$ enable step $B$?''), which is the primitive our project needs.

\paragraph{Listwise.} A whole candidate list enters one prompt and the model
emits a permutation using global context. RankGPT \citep{rl:sun2023rankgpt}
established the sliding-window protocol (window 20, step 10 over a top-100 list,
candidates labelled \texttt{[1] [2]\ldots} and reordered as \texttt{[2] > [3] >
[1]}), and distilled GPT into open students --- notably a 435M DeBERTa-v3-large
student trained with a RankNet pairwise loss that reportedly surpassed the prior
monoT5-3B state of the art on BEIR. RankVicuna \citep{rl:pradeep2023rankvicuna}
and RankZephyr \citep{rl:pradeep2023rankzephyr} are open 7B listwise rerankers
distilled from GPT-3.5/4 permutations; RankZephyr's two-stage distillation
(${\sim}100$K RankGPT-3.5 then ${\sim}5$K RankGPT-4 permutations) matches or
exceeds RankGPT-4 and is robust to input order. Listwise is cheap per list (few
calls) but suffers strong \emph{position bias} and finite context.

\paragraph{Setwise.} A set of $c{+}1$ documents enters each call and the model
selects only the most relevant; plugged into heapsort/bubblesort this cuts sort
complexity to $O(k\log_c N)$ \citep{rl:zhuang2024setwise}. On Flan-T5-large/DL19,
setwise heapsort used roughly half the inferences of pairwise heapsort at
comparable effectiveness. The ``pick the most relevant of a set'' objective is
the unit on which Rank-R1 later builds its RL.

\paragraph{Injecting side signals.} InsertRank
\citep{rl:chordia2025insertrank} appends each candidate's BM25 lexical score
into a listwise prompt to ground reasoning on reasoning-intensive queries (BRIGHT,
R2MED), giving consistent NDCG@10 gains with no fine-tuning. The lesson is that a
cheap structural prior injected into the prompt measurably helps --- analogous to
injecting a bi-encoder similarity or a workflow-position prior into a causal
reranker.

\subsection{Reasoning rerankers}
\label{sec:reasoning-rerankers}
Reasoning rerankers apply the DeepSeek-R1 recipe \citep{rl:deepseekr1} ---
\texttt{<think>}/\texttt{<answer>} formatting plus a verifiable relevance reward
optimised by GRPO --- to ranking.

\paragraph{Rank-R1.}
\citet{rl:zhuang2025rankr1} take the Setwise prompt, add a reasoning
instruction, and train the reasoning ability \emph{purely with RL} on relevance
labels (no reasoning supervision). On Qwen2.5 at 3B/7B/14B, GRPO uses a binary
reward (1 iff the output matches the \texttt{<think>}/\texttt{<answer>} format
\emph{and} the selected label is the relevant document) and matches the Setwise
SFT baseline using only ${\sim}18\%$ of its data, with the largest gains on
reasoning-intensive (BRIGHT) queries.

\paragraph{REARANK.}
\citet{rl:zhang2025rearank} train a \emph{listwise} reasoning agent
(REARANK-7B, on Qwen2.5-7B) with GRPO from only ${\sim}179$ annotated MS~MARCO
queries (expanded by augmentation to ${\sim}12$K instances). Its composite
reward $r = 0.8\,r_{\text{rank}} + 0.1\,r_{\text{fmt1}} + 0.1\,r_{\text{fmt2}}$
uses a min--max-normalised NDCG@10 \emph{improvement}
$(\mathrm{NDCG}_{\text{rerank}}-\mathrm{NDCG}_{\text{init}})/
(\mathrm{NDCG}_{\text{opt}}-\mathrm{NDCG}_{\text{init}})$ plus two format terms.
It matches GPT-4 on TREC DL19/20 and surpasses it on BRIGHT, demonstrating that a
listwise reasoning reranker --- whose permutation output is exactly a step
ordering --- can be learned from a tiny label set.

\paragraph{Rank-GRPO.}
\citet{rl:rankgrpo2025} generalise GRPO for ranked outputs by treating each
\emph{rank position} (not each token or sequence) as the unit of advantage
estimation and importance weighting, removing non-causal credit assignment and
the quality decay toward the end of a generated list. Although framed for
conversational recommendation, the rank-level credit assignment is directly
relevant to any listwise reasoning reranker.

\paragraph{When chain-of-thought \emph{hurts}: the think-free fix.} A central,
initially counter-intuitive finding is that naive chain-of-thought (CoT) can
\emph{degrade} pointwise rerankers \citep{rl:cotfallsshort2025}. The mechanism:
ranking needs a precise, calibrated \emph{scalar} comparable across candidates,
whereas free-form CoT (i) injects variance into the verbalised score, (ii) lets
the model ``overthink'' and drift from the query, and (iii) consumes the output
distribution that a clean score token would otherwise occupy --- all at large
latency cost. TFRank \citep{rl:tfrank2025} resolves this with a
\emph{think-mode switch}: the model is \emph{trained} with explicit CoT
(internalising reasoning into its weights) but runs \emph{think-free} at
inference, emitting only a pointwise score with no reasoning chain. Built on the
Qwen3 series with released checkpoints at \textbf{0.6B}, 1.7B, 4B and 8B, a
\emph{1.7B} TFRank competes with 7B reasoning rerankers on BRIGHT at ${\sim}4\times$
fewer parameters and an order-of-magnitude faster inference. Its score is an
expected-value over a five-level ordinal scale (0--4) read from the top-20 logit
tokens, with a neutral 0.5 fallback --- precisely the role a cross-encoder
classifier plays, but reasoning-distilled. The TFRank \textbf{0.6B} checkpoint is
the existence proof that a sub-1B think-free reasoning scorer is achievable rather
than speculative. Rank1 \citep{rl:weller2025rank1} occupies the opposite extreme
(emit CoT at inference as test-time compute), underscoring the
reason-then-rank{} versus reason-then-score{} latency trade-off.

\subsection{Reinforcement learning for ranking and reasoning at sub-1B scale}
\label{sec:rl-sub1b}
\paragraph{GRPO and its core mechanism.} Group Relative Policy Optimization
\citep{rl:shao2024grpo}, scaled in DeepSeek-R1 \citep{rl:deepseekr1}, removes
PPO's value network. For a prompt $q$ it samples a group of $G$ completions,
scores each with a reward $r_i$, and assigns every token of completion $i$ the
group-relative advantage
\begin{equation}
A_i = \frac{r_i - \mathrm{mean}(r_1,\ldots,r_G)}{\mathrm{std}(r_1,\ldots,r_G)},
\label{eq:grpo-adv}
\end{equation}
optimised with PPO's clipped surrogate plus a KL penalty to a frozen reference.
The group mean replaces the learned critic, halving memory --- attractive on
constrained GPUs --- but rollout cost scales with $G$, and groups that are
all-correct or all-wrong yield zero advantage.

\paragraph{GRPO variants.} \emph{Dr.\ GRPO} \citep{rl:liu2025drgrpo} identifies
two optimisation biases in Eq.~\ref{eq:grpo-adv}: response-length normalisation
(which inflates length, especially of wrong answers) and per-question std
normalisation (which over-weights easy/hard prompts). It normalises by a fixed
constant and drops the std term, cutting incorrect-response length by ${\sim}38\%$
while preserving accuracy. \emph{DAPO} \citep{rl:yu2025dapo} adds Clip-Higher
(decoupled clip range $\varepsilon_{\text{low}}{=}0.2$,
$\varepsilon_{\text{high}}{=}0.28$, fixing entropy collapse), dynamic sampling
(discarding zero-gradient groups), token-level loss, and overlong-reward shaping;
it removes the KL term and reaches AIME-2024$=50$ on Qwen2.5-32B in half the
training steps. \emph{S-GRPO} \citep{rl:sgrpo2025} samples serial early-exit
positions on a single chain with decaying reward, cutting sequence length 35--61\%
with no accuracy loss --- useful when a binary causal decision should not invite
runaway CoT.

\paragraph{PPO versus DPO versus GRPO.} PPO \citep{rl:schulman2017ppo} is
actor--critic: a learned value baseline is lower-variance but unstable and
doubles memory; it can edge out DPO on reasoning in controlled studies
\citep{rl:ivison2024ppodpo}. DPO \citep{rl:rafailov2023dpo} skips the RL loop and
reward model, optimising a classification loss directly on (chosen, rejected)
pairs --- cheap and stable, but it does not explore. GRPO is the middle ground for
\emph{verifiable} reasoning (no reward model, no critic). Because the causal
direction/order ground truth is verifiable, GRPO/Dr.\ GRPO is the natural fit;
DPO over (correct order $\succ$ corrupted order) pairs is a cheaper baseline.

\paragraph{Reward design.} The de-facto pattern is a sum of independent terms
\citep{rl:deepseekr1}: a \emph{format} reward ($+1$ if the output matches the
required schema) plus a \emph{correctness} reward checked by a deterministic
verifier. For binary causal direction the verifier is trivial
(predicted$=$gold). For \emph{ordering}, sparse exact-match wastes gradient, so
order-sensitive rewards are preferred: REARANK's NDCG-improvement
\citep{rl:zhang2025rearank}, or a Kendall-$\tau$/Spearman correlation between the
predicted and gold permutation (dense in $[-1,1]$), optionally plus an
exact-permutation bonus and a format term. Dense, order-aware rewards matter at
small scale because they avoid the all-wrong degenerate groups that give zero
advantage.

\paragraph{Distill-then-RL, not RL-from-scratch.} The strongest evidence for the
sub-1B regime is that \emph{distillation (SFT on reasoning traces) should precede
any RL}. DeepSeek-R1 \citep{rl:deepseekr1} reports that distilling a large model
into a small dense model beats large-scale RL run directly on that small model,
which ``requires enormous computational power and may not even achieve the
performance of distillation.'' Open-RS \citep{rl:dang2025openrs} confirms the
nuance: GRPO on 7K math samples lifted an \emph{already-distilled} 1.5B model
cheaply (${\sim}\$42$, 4$\times$A40), but observed degradation after 150--200
steps, KL instability and length truncation --- small-model RL is fragile. Tina
\citep{rl:wang2025tina} reproduced a competitive 1.5B reasoning model with
LoRA-GRPO for ${\sim}\$9$. ``SFT Memorizes, RL Generalizes''
\citep{rl:chu2025sftrl} gives the canonical justification: RL with outcome reward
generalises better out of distribution, \emph{but SFT is still required first} to
stabilise output format so RL pays off. A weak $<\!1$B base lacks the latent
reasoning for RL to discover much unaided; seed it with SFT, then run a short,
well-regularised GRPO polish (tens-to-low-hundreds of steps).

\subsection{Distilling reasoning into small / non-generative models}
\label{sec:distillation}
\paragraph{Rationale-as-auxiliary-objective.} Distilling Step-by-Step
\citep{rl:hsieh2023distillstep} trains a student with \emph{two} objectives ---
predict the label, and (auxiliary) produce the teacher's rationale --- so the
rationale is a \emph{training target, not an inference input}. A 770M T5 thereby
beat a 540B PaLM few-shot teacher using 80\% of the labelled data, and matched
full-data fine-tuning with as little as 12.5\% on e-SNLI. This is the conceptual
ancestor of internalised reasoning, and the precedent for an explanation-aligned
auxiliary head on a discriminative scorer; e-SNLI \citep{rl:camburu2018esnli} is
the canonical dataset pairing labels with natural-language explanations for
exactly this predict-and-explain supervision.

\paragraph{Distillation versus on-small-model RL.} R1-Distill
\citep{rl:deepseekr1} curated ${\sim}800$K R1 traces and applied \emph{pure SFT}
to open bases (Qwen2.5 1.5B--32B, Llama 8B/70B); the distilled 32B (AIME-2024
72.6 / MATH-500 94.3) far exceeds RL-from-scratch on the same base (47.0 / 91.6),
the cleanest published confirmation of the distill-first verdict. SCoTD
\citep{rl:li2023scotd} shows 125M--1.3B students gain CoT ability when trained on
many diverse teacher chains, and Magister et al.\ \citep{rl:magister2022teaching}
established that CoT fine-tuning transfers reasoning to small students.

\paragraph{Label-only rationale bootstrapping.} When only task labels exist (the
project's situation --- causal-direction labels, no gold rationales), rationales
can be self-generated. STaR \citep{rl:zelikman2022star} keeps self-produced
rationales that reach the correct answer and \emph{rationalises} the rest by
conditioning on the gold answer; rejection-sampling fine-tuning (RFT)
\citep{rl:yuan2023rft} is its one-shot best-of-$N$ variant. COLLATE
\citep{rl:patnaik2025collate} goes further with a teacher-free signal: it ranks
candidate rationales by how much each \emph{raises the model's own likelihood of
the gold answer}, forming DPO preference pairs --- no judge, no larger teacher,
needing only labels, and keeping a sample only when the winning rationale truly
raises gold-answer likelihood.

\paragraph{Rationale distillation into a cross-encoder (think-free, zero added
latency).} The unifying recipe is to add a rationale-derived auxiliary loss to a
discriminative head: a generative \texttt{[rationale]} decoding head (Distilling
Step-by-Step), or an embedding-alignment head that pulls the classifier's hidden
state toward a (frozen) rationale embedding via cosine/contrastive loss. TFRank's
think-free score \citep{rl:tfrank2025} is the ranking instantiation. At inference
only the label/score head runs, so reasoning improves accuracy at \emph{zero}
added latency --- exactly what a production DeBERTa-class cross-encoder needs,
since it cannot emit \texttt{<think>} tokens but can be regularised toward
reasoning-relevant features.

\subsection{Relevance to causal embedding for AEP/AJO}
This section is the backbone of the project's core thread: \emph{add reasoning to
a $<\!1$B causal scorer}. The repository's strongest ordering model is a
DeBERTa-v3-large cross-encoder, which cannot generate reasoning tokens, so genuine
reason-then-answer must be \emph{internalised}: the literature converges on
TFRank's think-free pointwise scorer (with a demonstrated 0.6B checkpoint) and
Distilling-Step-by-Step / e-SNLI-style auxiliary rationale heads as the way to
get reasoning gains at the latency of the current cross-encoder. The training
verdict is equally clear: with ${\sim}100$K--260K verifiable labels but no gold
rationales, the right pipeline is \emph{bootstrap rationales} (STaR/RFT/COLLATE)
$\rightarrow$ \emph{SFT-distil} the behaviour into a sub-1B (or cross-encoder)
scorer $\rightarrow$ a short \emph{GRPO/Dr.\ GRPO} polish with a format reward plus
a verifiable correctness reward for the pairwise causal direction and a dense
NDCG/Kendall-$\tau$ reward for the ordering head --- never RL from scratch. The
pairwise causal-precedence judgments this scorer produces are the input to the
global-ordering machinery of Sec.~\ref{sec:ordering}.
